\chapter{Research planning and doctoral education}
\label{chap:research-planning}

This chapter establishes the operational roadmap for the doctoral project, bridging the theoretical problem formulation presented in \cref{chap:research-gap-analysis} and the technical feasibility baselines established in \cref{chap:feasibility-study} with a structured four-year execution strategy.
To guarantee both scientific rigour and formal academic progression, the planning is structured into two main tracks:
\cref{sec:publication-plan} outlines the scientific dissemination plan, mapping the core research questions to discrete publishable topics across premier journals and conferences, alongside a top-down project execution timeline.
Subsequently, \cref{sec:doctoral-education-plan} details the formal doctoral education plan formulated in accordance with the regulations of the EEMCS Graduate School, demonstrating the fulfilment of transferable, research, and discipline-specific credit requirements.


\section{Publication plan and dissemination strategy}
\label{sec:publication-plan}

The primary mechanism for validating and disseminating the contributions of this doctoral work is the publication of peer-reviewed conference proceedings and journal articles.
Following the research gaps identified in \cref{chap:literature-review} and the problem formulation in \cref{sec:research-questions}, the overarching doctoral research is partitioned into ten publishable topics categorized under the four central research questions.


\subsection{Research-question breakdown and candidate venues}
\label{sec:rq-breakdown-venues}

To maximize scientific impact and maintain alignment with the targeted domain communities (antennas, radar signal processing, remote sensing, and intelligent transportation), candidate dissemination venues have been selected based on their peer-review rigor and domain relevance.

\paragraph{Research question 1.}
This research epoch focuses on whether integrating a high-resolution architecture with Doppler-resolved polarimetric processing can significantly enhance the classification accuracy of dynamic \glspl{vru}.
\begin{itemize}
    \item \textbf{Topic 1.1:} Design an X-band prototype of a large-aperture MIMO array with enhanced polarization purity.
    \begin{itemize}
        \item \textit{Focus:} Validating the baseline planar aperture-coupled series patch topology and stripline feeding network designed in \cref{sec:preliminary-antenna-design}.
        \item \textit{Candidate venues:} European Conference on Antennas and Propagation (EuCAP), European Microwave Week (EuMW), IEEE International Symposium on Antennas and Propagation (IEEE ISAP).
    \end{itemize}
        
    \item \textbf{Topic 1.2:} Design a first-iteration system leveraging mature technology platforms and uniform array synthesis.
    \begin{itemize}
        \item \textit{Focus:} Extends the validated X-band baseline to a \qty{79}{\giga\hertz} \gls{siw}/planar demonstrator, establishing a uniform dual-polarized array for initial radar measurement campaigns.
        \item \textit{Candidate venues:} IEEE Transactions on Antennas and Propagation (IEEE TAP), IEEE Antennas and Wireless Propagation Letters (IEEE AWPL), EuCAP.
    \end{itemize}
    
    \item \textbf{Topic 1.3:} Design a second-iteration system leveraging acquired experience and measurement campaign insights.
    \begin{itemize}
        \item \textit{Focus:} Refines the front-end architecture by integrating optimized polarimetric topologies and calibrated feed networks derived from initial experimental campaigns.
        \item \textit{Candidate venues:} IEEE TAP, IEEE Transactions on Microwave Theory and Techniques (IEEE T-MTT).
    \end{itemize}
\end{itemize}

\paragraph{Research question 2.}
This epoch investigates which polarimetric decomposition methods are most robust under non-stationary, grazing-incidence, and multipath automotive conditions, adapting them to exploit Doppler-resolved features.
\begin{itemize}
    \item \textbf{Topic 2.1:} Develop a polarimetric MIMO simulation platform linking full-wave array models with SBR+ ray tracing.
    \begin{itemize}
        \item \textit{Focus:} Establishes the hybrid simulation framework in Ansys Electronics Desktop presented in \cref{sec:simulation-platform}.
        \item \textit{Candidate venues:} IEEE RadarConf, EuMW / EuRAD, International Radar Symposium (IRS).
    \end{itemize}
    
    \item \textbf{Topic 2.2:} Perform perturbation analysis of polarimetric decomposition methods applicable to automotive imaging.
    \begin{itemize}
        \item \textit{Focus:} Maps algorithm breakdown thresholds against amplitude imbalances, phase errors, and cross-polar leakage injected into scattering matrices.
        \item \textit{Candidate venues:} IEEE Transactions on Radar Systems (IEEE T-RS), International Geoscience and Remote Sensing Symposium (IEEE IGARSS).
    \end{itemize}
    
    \item \textbf{Topic 2.3:} Adapt polarimetric processing algorithms to exploit Doppler-resolved features for VRU classification.
    \begin{itemize}
        \item \textit{Focus:} Formulates dynamic polarimetric-Doppler feature maps and validates classification performance on empirical pedestrian and cyclist radar data.
        \item \textit{Candidate venues:} IEEE Transactions on Geoscience and Remote Sensing (IEEE TGRS), IEEE IGARSS, IEEE Transactions on Intelligent Transportation Systems (IEEE T-ITS), IEEE Intelligent Vehicles Symposium (IEEE IV).
    \end{itemize}
\end{itemize}

\paragraph{Research question 3.}
This epoch links algorithm sensitivity directly to front-end electromagnetic figures of merit.
\begin{itemize}
    \item \textbf{Topic 3.1}: Link antenna figures of merit to sensitivity parameters of polarimetric decomposition methods.
    \begin{itemize}
        \item \textit{Focus:} Quantifies the mathematical mapping between front-end non-idealities (\gls{xpd}, port isolation, squint) and decomposition parameter variances.
        \item \textit{Candidate venues:} IEEE T-RS, EuCAP, EuMW, IEEE ISAP.
    \end{itemize}
    
    \item \textbf{Topic 3.2}: Derive fundamental lower bounds on polarimetric antenna requirements based on sensitivity analyses.
    \begin{itemize}
        \item \textit{Focus:} Formulates explicit minimum hardware bounds (e.g. minimum allowable \gls{xpd} across scan angle) required to prevent classification breakdown.
        \item \textit{Candidate venues:} IEEE TGRS, IEEE T-RS.
    \end{itemize}
\end{itemize}

\paragraph{Research question 4.}
This epoch addresses how \gls{mimo} virtual array manifolds can be optimized specifically for polarimetric fidelity to mitigate phase centre displacement and cross-polar coupling.
\begin{itemize}
    \item \textbf{Topic 4.1}: Explore the effects of bistatic phase centre displacement of antenna pairs on polarimetric phase distortion.
    \begin{itemize}
        \item \textit{Focus:} Analyses virtual channel parallax and angle-dependent phase squint in interleaved single-polarized and dual-polarized \gls{mimo} arrays.
        \item \textit{Candidate venues:} EuCAP, IEEE RadarConf, IRS, IEEE ISAP.
    \end{itemize}
    
    \item \textbf{Topic 4.2}: Construct objective functions accounting for polarimetric distortion for 'polarimetry-first' MIMO layouts.
    \begin{itemize}
        \item \textit{Focus:} Develops mathematical optimization frameworks for virtual array synthesis that minimize phase centre disparity across co- and cross-polar channels.
        \item \textit{Candidate venues:} IEEE AWPL, IEEE T-RS, EuMW / EuRAD, IEEE RadarConference
    \end{itemize}
\end{itemize}

The complete mapping of research questions, publishable topics, target publication types, candidate venues, and temporal milestones is presented in the landscape overview on the following page (\cref{tab:publication-venues-summary,fig:publication-gantt}).

\begin{landscape}
\begin{table}[!p]
    \centering
    \caption{Summary of research questions, publishable topics, and candidate dissemination venues.}
    \label{tab:publication-venues-summary}
    \small
    \begin{tabularx}{\linewidth}{lXXl}
        \toprule
        \textbf{Topic} & \textbf{Focus / Title} & \textbf{Target output} & \textbf{Candidate venues} \\
        \midrule
        Topic 1.1 & X-band prototype array & Conference paper & EuCAP, EuMW, IEEE ISAP \\
        Topic 1.2 & 1st iteration \qty{79}{GHz} system & Letter / journal & IEEE AWPL, IEEE TAP, EuCAP \\
        Topic 1.3 & 2nd iteration refined system & Journal paper & IEEE TAP, IEEE T-MTT \\
        \midrule
        Topic 2.1 & Hybrid SBR+ simulation platform & Conference paper & IEEE RadarConf, EuMW, IRS \\
        Topic 2.2 & Decomposition perturbation analysis & Journal / Conference & IEEE T-RS, IEEE IGARSS \\
        Topic 2.3 & Doppler-resolved VRU classification & Journal paper & IEEE TGRS, IEEE IGARSS, IEEE T-ITS, IEEE IV \\
        \midrule
        Topic 3.1 & Antenna metrics vs sensitivity & Conference / Journal & IEEE T-RS, EuCAP, EuMW, IEEE ISAP \\
        Topic 3.2 & Fundamental hardware lower bounds & Journal paper & IEEE TGRS, IEEE T-RS \\
        \midrule
        Topic 4.1 & Phase centre displacement distortion & Conference paper & EuCAP, IEEE RadarConf, IRS, IEEE ISAP \\
        Topic 4.2 & Polarimetry-first topology synthesis & Letter / Journal & IEEE AWPL, IEEE T-RS, EuMW, IEEE RadarConf \\
        \bottomrule
    \end{tabularx}
\end{table}

\vfill

\begin{figure}[!p]
    \centering
    \small
    \begin{ganttchart}[
        x unit=0.42cm,
        y unit title=0.55cm,
        y unit chart=0.45cm,
        vgrid={*1{black!15}},
        hgrid,
        time slot format=isodate,
        time slot unit=month,
        title/.style={fill=black!10, draw=black!30, font=\small\bfseries},
        title label font=\small\bfseries,
        bar/.style={fill=blue!30, draw=blue!70},
        bar progress label font=\scriptsize,
        milestone/.style={fill=red!70, draw=red!90}
    ]{2025-09-01}{2029-08-31}
        \gantttitlecalendar{year} \\
        \gantttitle{}{1} \gantttitle{Q4}{3}
        \gantttitle{Q1}{3} \gantttitle{Q2}{3} \gantttitle{Q3}{3} \gantttitle{Q4}{3}
        \gantttitle{Q1}{3} \gantttitle{Q2}{3} \gantttitle{Q3}{3} \gantttitle{Q4}{3}
        \gantttitle{Q1}{3} \gantttitle{Q2}{3} \gantttitle{Q3}{3} \gantttitle{Q4}{3}
        \gantttitle{Q1}{3} \gantttitle{Q2}{3} \gantttitle{}{2} \\
        % Intro / Background (Gray)
        \ganttbar[bar/.style={fill=gray!30, draw=gray!70}]{Onboarding}{2025-09-01}{2025-09-30} \\
        \ganttbar[bar/.style={fill=gray!30, draw=gray!70}]{Lit. review}{2025-09-01}{2025-12-31} \\
        % RQ1: Integration & VRU (Blue family)
        \ganttbar[bar/.style={fill=blue!30, draw=blue!75}]{Topic 1.1}{2026-01-01}{2026-10-31} \\
        \ganttbar[bar/.style={fill=blue!70!black, draw=blue!90!black}]{Topic 1.2}{2026-11-01}{2027-03-31} \\
        \ganttbar[bar/.style={fill=blue!30, draw=blue!75}]{Tape-out 1}{2027-04-01}{2027-08-31} \\

        \ganttbar[bar/.style={fill=blue!70!black, draw=blue!90!black}]{Topic 1.3}{2028-02-01}{2028-03-31} \\
        \ganttbar[bar/.style={fill=blue!30, draw=blue!75}]{Tape-out 2}{2028-04-01}{2028-10-31} \\
        % RQ2: Robust decomposition & Doppler (Teal)
        \ganttbar[bar/.style={fill=teal!35, draw=teal!75}]{Topic 2.1}{2026-06-01}{2026-12-31} \\
        \ganttbar[bar/.style={fill=teal!35, draw=teal!75}]{Topic 2.2}{2027-01-01}{2027-04-30} \\
        \ganttbar[bar/.style={fill=teal!35, draw=teal!75}]{Topic 2.3}{2028-05-01}{2028-11-30} \\
        % RQ3: Fundamental lower bounds (Orange)
        \ganttbar[bar/.style={fill=orange!35, draw=orange!80}]{Topic 3.1}{2027-05-01}{2027-08-31} \\
        \ganttbar[bar/.style={fill=orange!35, draw=orange!80}]{Topic 3.2}{2027-09-01}{2027-12-31} \\
        % RQ4: Topology optimization (Violet)
        \ganttbar[bar/.style={fill=violet!30, draw=violet!75}]{Topic 4.1}{2027-04-01}{2027-07-31} \\
        \ganttbar[bar/.style={fill=violet!30, draw=violet!75}]{Topic 4.2}{2027-08-01}{2028-01-31} \\
        % Thesis (Dark Blue / Navy)
        \ganttbar[bar/.style={fill=blue!85!black, draw=black}]{Thesis}{2029-03-01}{2029-08-31}
    \end{ganttchart}
    \caption{Four-year PhD publication and project execution timeline (colour-coded by research question).}
    \label{fig:publication-gantt}
\end{figure}
\end{landscape}


\section{Doctoral education plan}
\label{sec:doctoral-education-plan}

In accordance with the regulations of the EEMCS Graduate School, doctoral candidates must accumulate a minimum of \qty{45}{ECTS} credits evenply split across three competence categories: transferable, research, and discipline-specific skills.
This education plan is structured to provide rigorous technical training during the early phases of the PhD while supporting career development and effective research communication throughout the trajectory.


\subsection{Graduate school credit requirements}
\label{sec:gs-credit-requirements}

The credit distribution across the three Graduate School brackets is summarized as follows:
\begin{enumerate}
    \item \textbf{Transferable skills.} Courses designed to develop personal effectiveness, project management, leadership, and career development. A minimum of \qty{1}{ECTS} must be dedicated to \emph{career development}.
    \item \textbf{Research skills.} Activities focusing on research methodology, scientific writing, information skills, and teaching responsibilities. A minimum of \qty{5}{ECTS} must be accumulated via Learning on the Job (LotJ) activities.
    \item \textbf{Discipline-specific skills.} Advanced academic courses, summer schools, and specialized workshops directly aimed at acquiring domain-specific knowledge.
\end{enumerate}


\subsection{Transferable skills}
\label{sec:transferable-skills}

\Cref{tab:transferable-skills} details the selection of activities within the \emph{transferable skills} category.
The mandatory career-development requirement is satisfied via module T4.G5 (\qty{1.0}{ECTS}).

\begin{table}[!ht]
    \centering
    \caption{Doctoral education plan: transferable skills.}
    \label{tab:transferable-skills}
    \small
    \begin{tabularx}{\textwidth}{Xccc}
        \toprule
        \textbf{Activity / Course code} & \textbf{Credits} & \textbf{Year} & \textbf{Status} \\
        \midrule
        T4.G1 - A \textbar{} PhD Start-up Module A & 1.5 & 1 & Completed \\
        T4.G1 - B \textbar{} PhD Start-up Module B & 0.5 & 1 & Completed \\
        T1.A9 \textbar{} Scientific Text Processing with \LaTeX & 1.0 & 1 & Completed \\
        T1.A7 \textbar{} Data Vizualisation & 1.0 & 1 & Completed \\
        T4.B3 \textbar{} Time Management II & 1.0 & 1 & Completed \\
        T4.B5 \textbar{} Project Management for PhD Candidates & 2.0 & 1 & Completed \\
        T2.B1 \textbar{} Effective Interaction within Research Team & 2.5 & 1 & Enrolled \\
        T4.B2 \textbar{} Time Management I & 1.0 & 2--3 & Planned \\
        T2.D2 \textbar{} Managing Myself, Leading Others & 1.5 & 2--3 & Planned \\
        T4.A2 \textbar{} Smart about Stress & 2.0 & 2--3 & Planned \\
        T4.G5 \textbar{} Career Development -- Personal Branding$^\ast$ & 1.0 & 3--4 & Planned \\
        \midrule
        \textbf{Total} & \textbf{15.0/15} & & \\
        \midrule
        \multicolumn{4}{l}{\footnotesize $^\ast$ Denotes course fulfilling mandatory Graduate School career-development requirement.} \\
        \bottomrule
    \end{tabularx}
\end{table}


\subsection{Research skills and learning on the job}
\label{sec:research-skills}

\Cref{tab:research-skills} outlines the planned activities within the \emph{research skills} bracket.
In total, \qty{14.0}{ECTS} are earned through LotJ activities, exceeding the Graduate School threshold of \qty{5.0}{ECTS}.

\begin{table}[!ht]
    \centering
    \caption{Doctoral education plan: research skills.}
    \label{tab:research-skills}
    \small
    \begin{tabularx}{\textwidth}{Xcccc}
        \toprule
        \textbf{Activity / Course code} & \textbf{Credits} & \textbf{LotJ} & \textbf{Year} & \textbf{Status} \\
        \midrule
        Teaching assistance: laboratory/tutorial support & 2.0 & Yes & 1 & Completed \\
        Teaching assistance: designing examination tasks & 1.0 & Yes & 1 & Completed \\
        Teaching assistance: technical support \& exam grading & 2.0 & Yes & 1 & Completed \\
        R1.A2 \textbar{} The Informed Researcher & 1.5 & No & 1 & Completed \\
        Writing a research proposal & 2.0 & Yes & 1 & In progress \\
        Writing a conference paper & 1.0 & Yes & 1--2 & Planned \\
        Writing journal articles ($2\times$) & 6.0 & Yes & 2--3 & Planned \\
        \midrule
        \textbf{Core Total} & \textbf{15.5/15} & \textbf{14.0/5} & & \\
        % \midrule
        % \multicolumn{5}{l}{\textit{Elective candidate activities:}} \\
        % International research internship ($\ge 1$~month) & 2.0 & Yes & 1--3 & Optional \\
        % R2.B2 \textbar{} Statistics for PhD Research & 4.0 & No & 1--3 & Optional \\
        % R2.B3 \textbar{} Data Visualization in Python & 1.0 & No & 1--3 & Optional \\
        \bottomrule
    \end{tabularx}
\end{table}


\subsection{Discipline-specific skills}
\label{sec:discipline-skills}

\Cref{tab:discipline-skills} details the specialized coursework selected to deepen theoretical knowledge in millimetre-wave antennas, phased array synthesis, and high-frequency electromagnetics.
Due to the intense workload associated with advanced technical modules, discipline skills are allocated primarily to the first two years of the doctoral trajectory.

\begin{table}[!ht]
    \centering
    \caption{Doctoral education plan: discipline skills.}
    \label{tab:discipline-skills}
    \small
    \begin{tabularx}{\textwidth}{Xccc}
        \toprule
        \textbf{Activity / Course title} & \textbf{Credits} & \textbf{Year} & \textbf{Status} \\
        \midrule
        International Summer School on Phased Array Systems & 5.0 & 1 (Q1) & Completed \\
        EE4016 \textbar{} Antenna Systems & 5.0 & 1 (Q3) & Completed \\
        Arrays and Reflectarrays \textbar{} ESoA 2026 Course & 5.0 & 2 (Sep 2026) & Planned \\
        \midrule
        \textbf{Total} & \textbf{15.0/15} & & \\
        \bottomrule
    \end{tabularx}
\end{table}
